\subsection{Experimental Setup}
We utilize the M3FD and FLIR marine datasets, splitting data by distinct weather and lighting conditions (e.g., train on Day/Clear, test on Night/Fog) to enforce domain generalization.

\subsection{Baselines}
We compare against five distinct strategies: (1) 	extbf{Classical Registration:} Fixed SIFT/ORB + RANSAC; (2) 	extbf{Fixed Augmentation:} Standard supervised learning with static random flip/crop; (3) 	extbf{RandAugment:} Randomized search without context; (4) 	extbf{AutoAugment:} Offline search for a fixed best policy; and (5) 	extbf{Context-Free RL:} An ablation of our method ignoring state inputs.

\subsection{Metrics}
Primary metrics include Mean Average Precision (mAP@50) for detection and Mean Reprojection Error (MRE) for alignment quality. Secondary metrics assess robustness to induced misalignment and sensor degradation (simulated thermal crossover).

\subsection{Hypothesis Validation}
Success is defined as the RL policy achieving higher mAP than AutoAugment on the hardest test split. We specifically analyze the "Action Distribution" to verify if the policy correctly switches strategies—for example, choosing robust rigid registration during high-noise night scenes where homography estimation might collapse (Fig.~
ef{fig:published-evidence}).
